\documentclass[../../main]{subfiles}
\begin{document}

This thesis presents a systematic investigation of the use of Large Language Models (LLMs)~\cite{LLMsIBM} within the domain of Constraint Programming (CP)~\cite{CPholygrail}, specifically in the task of algorithm selection~\cite{AlgSelProb}, referred to as ``solver selection'' in the context of NP-complete problems solved by constraint solvers. The objective is to assess whether general-purpose language models can effectively support or enhance decision processes that have traditionally relied on supervised learning.

Algorithm selection aims to identify the most suitable algorithm for solving a specific problem instance prior to execution. This issue arises in many domains because no single algorithm performs optimally across all problem instances. Traditional algorithm selection and portfolio construction methods typically treat the problem as a supervised classification or regression task~\cite{algSelRanking}.

LLMs are statistical language models that estimate probability distributions over token sequences. In practice, these distributions are parameterized by deep neural network architectures trained on large text corpora. They result from sustained advances in natural language processing and machine learning, and have recently demonstrated strong performance across a wide range of reasoning and generation tasks. Despite their broad adoption in general-purpose applications, their role in specialized technical domains such as CP, and specifically in automated solver selection, remains largely unexplored.

Constraint Programming is a declarative paradigm for solving combinatorial problems, particularly those arising in planning, scheduling, and resource allocation~\cite{CPholygrail}. Problems are modeled in terms of variables, domains, and constraints, and are processed by solvers, i.e., software systems that search for assignments satisfying the constraints and, in optimization settings, improving a given objective function. Different solvers rely on distinct underlying technologies and heuristics, which leads to performance variability across problem classes. Selecting an appropriate solver for a given instance is therefore a central challenge~\cite{evaluationMetaSolvers}.

This thesis investigates the problem of algorithm selection by using LLMs as decision-making components within a solver portfolio. This idea is inspired by~\cite{fromPFStoAS} and builds on portfolio-based approaches, where multiple solvers are available and a central strategy determines which ones to execute for a given instance. Traditional portfolio methods rely on complex supervised learning algorithms. In contrast, this work explores whether general-purpose LLMs, given suitable contextual information, can approximate or surpass these approaches without task-specific retraining.

Although the study is exploratory in scope, its experimental design is organized around a set of explicit working assumptions that motivate the progression of experiments. It is hypothesized that structured representations of problem instances will be more informative for solver selection than raw scripts. In addition, feature-based numerical summaries capturing structural properties of each instance are expected to provide stronger discriminative signals than text alone. Furthermore, the inclusion of solver descriptions in the prompt is expected to reduce biases arising from solver name recognition, enabling the model to reason on the basis of capability rather than prior association. Finally, sampling temperature is treated as a control variable regulating the trade-off between exploration and exploitation in solver choice, with lower values expected to favor stable, high-probability selections and higher values enabling broader exploration of the solver space.

The experimental evaluation relies on benchmark instances from the 2025 MiniZinc Challenge~\cite{PhilMznChallenge,MznResults25}. A set of general-purpose LLMs was selected for evaluation and conducted iterative experiments to assess whether these models could outperform existing single solvers, systematically testing different prompt configurations and progressively enriching the contextual information provided.

Initial experiments revealed substantive limitations in the ability of LLMs to support this task, producing only marginal improvements relative to individual baseline solvers. To mitigate these limitations, structured representations of the problem instances were incorporated through the feature extraction tool \texttt{mzn2feat}~\cite{mzn2feat}, which derives a set of static and dynamic numerical features from each instance, such as the number of variables and constraints. While this representation yielded the highest performance observed in our experiments, the overall improvement remained limited. Furthermore, prompts constructed from complete model-and-data encodings or from extensive feature vectors may become prohibitively large with respect to practical context-window constraints. This consideration motivated the development of \texttt{fzn2nl}~\cite{fzn2nl}, a tool designed to translate problem code into compact natural language descriptions, which were subsequently employed in additional experiments.

Overall, the results indicate that LLM-based solver selection is promising but fragile: the best configurations achieve only marginal gains over the single best solver in the full portfolio, and performance decreases in the restricted portfolio where finer discrimination is required. The \texttt{fzn2nl} representation shows mixed effects, not improving top-1 selection over the strongest feature-based variants but improving diversity and achieving the best parallel score in the restricted setting. The reasoning analysis further shows that LLMs can produce coherent solver-level and problem-level explanations, yet their final choices are often driven by prior associations and do not consistently align with empirical performance.

Beyond the empirical findings, this thesis contributes a reproducible evaluation protocol for LLM-based solver selection, a structured comparison of input representations, the deterministic \texttt{fzn2nl} pipeline for compact natural-language abstractions of FlatZinc models, and an interpretative framework based on controlled perturbation tests (including description swapping and solver anonymization) to analyze decision patterns in addition to predictive accuracy.

The structure of the thesis is as follows. \Cref{chap:techBg} provides the technical background required for the remainder of the work. \Cref{chap:methodology} describes the experimental methodology, evaluation metrics, and design choices underlying the study. \Cref{chap:expEval} presents the core experimental results, beginning with baseline tests based on minimal problem representations and progressively enriching the context with textual descriptions, structured features extracted from problem instances, and finally temperature tuning. \Cref{chap:parser} introduces a complementary tool developed during this research, \texttt{fzn2nl}, which translates FlatZinc models into natural language descriptions, and evaluates LLM performance when operating directly on such generated representations. \Cref{chap:reasoning} describes an experimental investigation aimed at understanding the level of comprehension that LLMs exhibit with respect to solvers and problem characteristics. \Cref{chap:conclusions} presents the conclusions.

\end{document}